\documentclass[12pt,a4paper]{article}
\usepackage[T1]{fontenc}
\usepackage{amsmath,amssymb,mathrsfs,tikz,times,pifont}
\usepackage{enumitem}
\usepackage{multicol}
\usepackage{lmodern}
\usetikzlibrary{trees}
\newcommand\circitem[1]{%
\tikz[baseline=(char.base)]{
\node[circle,draw=gray, fill=red!55,
minimum size=1.2em,inner sep=0] (char) {#1};}}
\newcommand\boxitem[1]{%
\tikz[baseline=(char.base)]{
\node[fill=cyan,
minimum size=1.2em,inner sep=0] (char) {#1};}}
\setlist[enumerate,1]{label=\protect\circitem{\arabic*}}
\setlist[enumerate,2]{label=\protect\boxitem{\alph*}}
%%%::::::by chnini ameur :::::::%%%
\everymath{\displaystyle}
\usepackage[left=1cm,right=1cm,top=1cm,bottom=1.7cm]{geometry}
\usepackage[colorlinks=true, linkcolor=blue, urlcolor=blue, citecolor=blue]{hyperref}
\usepackage{array,multirow}
\usepackage[most]{tcolorbox}
\usepackage{varwidth}
\usepackage{float} %pour utiliser l'option [H] qui force l'image à apparaître exactement à l'endroit où elle est placée dans le code.
\tcbuselibrary{skins,hooks}
\usetikzlibrary{patterns}
%%%::::::by chnini ameur :::::::%%%
\newtcolorbox{exa}[2][]{enhanced,breakable,before skip=2mm,after skip=5mm,
colback=yellow!20!white,colframe=black!20!blue,boxrule=0.5mm,
attach boxed title to top left ={xshift=0.6cm,yshift*=1mm-\tcboxedtitleheight},
fonttitle=\bfseries,
title={#2},#1,
% varwidth boxed title*=-3cm,
boxed title style={frame code={
\path[fill=tcbcolback!30!black]
([yshift=-1mm,xshift=-1mm]frame.north west)
arc[start angle=0,end angle=180,radius=1mm]
([yshift=-1mm,xshift=1mm]frame.north east)
arc[start angle=180,end angle=0,radius=1mm];
\path[left color=tcbcolback!60!black,right color = tcbcolback!60!black,
middle color = tcbcolback!80!black]
([xshift=-2mm]frame.north west) -- ([xshift=2mm]frame.north east)
[rounded corners=1mm]-- ([xshift=1mm,yshift=-1mm]frame.north east)
-- (frame.south east) -- (frame.south west)
-- ([xshift=-1mm,yshift=-1mm]frame.north west)
[sharp corners]-- cycle;
},interior engine=empty,
},interior style={top color=yellow!5}}
%%%%%%%%%%%%%%%%%%%%%%%
\usepackage{fancyhdr}
\usepackage{eso-pic}         % Pour ajouter des éléments en arrière-plan
% Commande pour ajouter du texte en arrière-plan
\usepackage{tkz-tab}
\AddToShipoutPicture{
    \AtTextCenter{%
        \makebox[0pt]{\rotatebox{80}{\textcolor[gray]{0.7}{\fontsize{5cm}{5cm}\selectfont PGB}}}
    }
}
\usepackage{lastpage}
\fancyhf{}
\pagestyle{fancy}
\renewcommand{\footrulewidth}{1pt}
\renewcommand{\headrulewidth}{0pt}
\renewcommand{\footruleskip}{10pt}
\fancyfoot[R]{
\color{blue}\ding{45}\ \textbf{2025}
}
\fancyfoot[L]{
\color{blue}\ding{45}\ \textbf{Prof:M. BA}
}
\cfoot{\bf
\thepage /
\pageref{LastPage}}
% Création du compteur pour les exercices
\newcounter{exercice}
\renewcommand{\theexercice}{\arabic{exercice}}  % Définit l'affichage du compteur en chiffres arabes

% Définir la commande \exo
\newcommand{\exo}{\refstepcounter{exercice}\textbf{Exercice \theexercice} }
\begin{document}
\renewcommand{\arraystretch}{1.5}
\renewcommand{\arrayrulewidth}{1.2pt}
\begin{tikzpicture}[overlay,remember picture]
    \node[draw=blue,line width=1.2pt,fill=purple,text=blue,inner sep=3mm,rounded corners,pattern=dots]at ([yshift=-2.5cm]current page.north) {\begingroup\setlength{\fboxsep}{0pt}\colorbox{white}{\begin{tabular}{|*1{>{\centering \arraybackslash}p{0.28\textwidth}} |*2{>{\centering \arraybackslash}p{0.2\textwidth}|} *1{>{\centering \arraybackslash}p{0.19\textwidth}|} }
                \hline
                \multicolumn{3}{|c|}{$\diamond$$\diamond$$\diamond$\ \textbf{Lycée de Dindéfélo}\ $\diamond$$\diamond$$\diamond$ } & \textbf{A.S. : 2026/2027}                                              \\ \hline
                \textbf{Matière: Mathématiques}                                                                                    & \textbf{Niveau : $1^{er}$}\textbf{L} & \textbf{Date: 05/11/2026} & \textbf{} \\ \hline
                \multicolumn{4}{|c|}{\parbox[c]{10cm}{\begin{center}
                                                                  \textbf{{\Large\sffamily Td Equations et Systèmes}}
                                                              \end{center}}}                                                                                                        \\ \hline
            \end{tabular}}\endgroup};
\end{tikzpicture}

\vspace{3cm}

\fbox{\textbf{\exo}}

Pour chacune des questions suivantes, une seule des réponses proposées est exacte.

\begin{enumerate}
  \item Le système 
  \[
  \begin{cases}
  2x + y - 5z = 0 \\
  \quad -10y + 2z = 20 \\
  \quad\quad\qquad 4z = 0
  \end{cases}
  \]
  a pour ensemble solution :
  \begin{align*}
  a)~S = \{(0,0,0)\} \qquad 
  b)~\varnothing \qquad 
  c)~S = \{(1,-2,0)\}
  \end{align*}

  \item Le système 
  \[
  \begin{cases}
  x + y + z = 15 \\
  -x - y + 2z = 0 \\
  2x + y - z = 8
  \end{cases}
  \]
  a pour ensemble solution :
  \begin{align*}
  a)~\{(9,1,5)\} \qquad 
  b)~\{(7,4,4)\} \qquad 
  c)~\{(3,7,5)\} \qquad 
  d)~\{(4,7,4)\}
  \end{align*}

  \item Le demi-plan d’inéquation \(2x - 5y + 9 < 0\) contient le point :
  \begin{align*}
  a)~(0,0) \qquad 
  b)~(-3,1) \qquad 
  c)~(-2,-1)
  \end{align*}

  \item Le système 
  \[
  \begin{cases}
  5x - 7y = -22 \\
  x + 2y < -1
  \end{cases}
  \]
  définit un ensemble qui contient le point :
  \begin{align*}
  a)~(1,1) \qquad 
  b)~(-3,1) \qquad 
  c)~(2,-5)
  \end{align*}
\end{enumerate}

\fbox{\textbf{\exo}}
Résoudre dans $\mathbb{R}^3$ les systèmes suivants, par la méthode du \emph{pivot de Gauss}.
\begin{multicols}{2}
\setlength{\columnseprule}{0.1mm} % La largeur de la ligne verticale entre les colonnes
\begin{enumerate}
  \item
  $
  \begin{cases}
  2x + y - 5z = 0\\[4pt]
  3x - y + 2z = -12\\[4pt]
  2x - 4y - 10z = 10
  \end{cases}
  $

  \item
  $
  \begin{cases}
  6x - 2y - 4z = 0\\[4pt]
  x + y - z = 1\\[4pt]
  5x - 5y + z = 1
  \end{cases}
  $

  \item
  $
  \begin{cases}
  3a + 2b - 5c = -8\\[4pt]
  -a + 4b + c = 6\\[4pt]
  2a - b + 3c = 5
  \end{cases}
  $

  \item (système à quatre inconnues $x,y,z,t$)
  \(
  \begin{cases}
  3x - 2y - 5z + t = -5\\[4pt]
  2x + 3y - 3z + 2t = -6\\[4pt]
  5x - 2y + 2z - 3t = -4\\[4pt]
  2x - 4y + z - t = 2
  \end{cases}
  \)
\end{enumerate}
\end{multicols}
\vspace{0.6cm}
\fbox{\textbf{\exo}}
À l'aide de graphiques différents, résoudre chacun des systèmes d'inéquations :
\begin{multicols}{2}
\setlength{\columnseprule}{0.1mm} % La largeur de la ligne verticale entre les colonnes
\begin{enumerate}
  \item
  $
  \begin{cases}
  3x + y < 2\\[4pt]
  2x - 3y > 1
  \end{cases}
  $

  \item
  $
  \begin{cases}
  x + y - 2 \le 0\\[4pt]
  x \ge 0\\[4pt]
  y < 3
  \end{cases}
  $

  \item
  $
  \begin{cases}
  4x - 5y + 6 > 0\\[4pt]
  5x + y - 7 < 0\\[4pt]
  x + 6y + 16 > 0
  \end{cases}
  $
\end{enumerate}
\end{multicols}
\fbox{\textbf{\exo}}
Dans un marché, trois enfants achètent les mêmes variétés de fruits, le 1er achète une orange, deux mandarines et une banane et paye 65 F; le 2ème achète deux oranges, une mandarine et trois bananes et paye 125 F et le 3ème achète trois oranges, une mandarine et une banane et paye 95 F. \\
Quel est le prix unitaire de chaque variété de fruit ?

\vspace{0.5cm}

\fbox{\textbf{\exo}}
Paul possède trois (03) sacs dont un (01) de maïs, un (01) d’igname et un (01) de riz. Les trois (03) sacs pèsent ensemble 120kg. La somme des poids du sac de maïs et du sac de riz est le double de celui du sac d’igname. S’il on ajoute 75kg au sac du riz, son poids sera le double de la somme des poids du sac du maïs et du sac d’igname. On désigne par le poids du sac de maïs, y celui du sac d’igname, et z celui du sac de riz

\begin{enumerate}
    \item Déterminer le système (S) que vérifient x, y et z
    \item En déduire le poids de chaque sac.
\end{enumerate}
\fbox{\textbf{\exo}}
\textbf{1) Résoudre dans $\mathbb{R}^3$ le système suivant, par la méthode du pivot de Gauss.}

\[
\text{a)} 
\begin{cases}
x + y + z = -1 \\
x - y + z = 2 \\
4x + 2y + z = -4
\end{cases}
\]

\textbf{2) Soit le polynôme $P(x) = x^3 + a x^2 + b x + c$ où $a$, $b$, $c$ sont trois nombres réels.}

Déterminer $a$, $b$, $c$ sachant que $P(1) = 0$, $P(-1) = 1$ et $P(2) = 4$.\\[0.3cm]
\textbf{3)} Factoriser $P(x)$.

\fbox{\textbf{\exo}}

Une entreprise fabrique deux types de produits, $A$ et $B$.

\begin{itemize}
    \item Chaque unité de produit $A$ nécessite 2 heures de travail en machine et 3 unités de matière première.
    \item Chaque unité de produit $B$ nécessite 4 heures de travail en machine et 2 unités de matière première.
\end{itemize}

L’entreprise dispose de :
\begin{itemize}
    \item 60 heures de travail en machine.
    \item 40 unités de matière première.
\end{itemize}

Les bénéfices pour chaque unité de produit $A$ et $B$ sont respectivement de 10 et 15.\\
Soit $x$ le nombre d’unités de $A$ et $y$ le nombre d’unités de $B$.

\begin{enumerate}
    \item Écrire l’expression du bénéfice à maximiser en fonction de $x$ et $y$.
    \item Écrire l’inéquation qui traduit les contraintes sur le nombre d’heures de machine.
    \item Écrire l’inéquation qui traduit les contraintes sur la quantité de matière première.
    \item Écrire le système des contraintes de l’entreprise.
    \item Déterminer combien d’unités de chaque produit doivent être fabriquées pour maximiser le bénéfice.
\end{enumerate}
\end{document}